% Part: computability
% Chapter: computability-theory
% Section: no-universal-function

\documentclass[../../../include/open-logic-section]{subfiles}

\begin{document}

\olfileid{cmp}{thy}{nou}
\olsection{सार्वत्रिक संगणनक्षम फलन नसते}

आंशिक संगणनक्षम फलनांसाठी सार्वत्रिक असे आंशिक संगणनक्षम फलन असले,
तरी सर्वत्र परिभाषित संगणनक्षम फलनांसाठी सार्वत्रिक असे सर्वत्र
परिभाषित संगणनक्षम फलन नसते.
%READERNOTE{OLCMP-015 — गोठवलेल्या इंग्रजीतील पहिल्या वाक्यात ``total for the partial computable functions'' असे आले आहे. मागील प्रमेय आणि पुढील विरोध पाहता येथे अभिप्रेत शब्द ``universal'' आहे; मराठीत तोच अर्थ घेतला आहे.}

\begin{thm}
\ollabel{thm:no-univ}
सार्वत्रिक संगणनक्षम फलन नसते. दुसऱ्या शब्दांत, असे कोणतेही
फलन $\fn{Un}'(k, x)$ घ्या की $f(x)$ हे सर्वत्र परिभाषित संगणनक्षम
फलन असल्यास, अशी नैसर्गिक संख्या~$k$ असते की प्रत्येक~$x$ साठी
$f(x) = \fn{Un}'(k,x)$; तर असे फलन संगणनक्षम नसते.
\end{thm}

\begin{proof}
सिद्धता साध्या विकर्णीकरणाने मिळते: $\fn{Un}'(k,x)$ हे सर्वत्र
परिभाषित आणि संगणनक्षम असेल, तर
\[
d(x) = \fn{Un}'(x, x) + 1
\]
हेही सर्वत्र परिभाषित आणि संगणनक्षम असेल. परंतु व्याख्येवरून
$d(k)$ हे $\fn{Un}'(k,k)$ च्या समान नाही. म्हणून प्रत्येक $k$
साठी $d(x)$ आणि~$\fn{Un}'(k, x)$ यांची मूल्ये किमान एका~$x$
साठी, म्हणजे $x = k$ साठी, वेगळी असतात.
\end{proof}

\begin{explain}
वरील \olref[uni]{thm:univ-comp} दाखवते की सार्वत्रिक फलन आंशिक
असण्याची परवानगी दिल्यास हे विकर्णीकरण टाळता येते. म्हणजे,
$\fn{Un}$ हे सर्वत्र परिभाषित संगणनक्षम फलनांसाठी सार्वत्रिक आहे;
पण ते स्वतः सर्वत्र परिभाषित नाही. आंशिक फलनांच्या बाबतीत
विकर्णीकरणाचा युक्तिवाद चालत नाही.
\end{explain}

\begin{prob}
  \olref{thm:no-univ} च्या सिद्धतेतील विकर्णीकरणाचा युक्तिवाद आंशिक
  फलनांच्या बाबतीत का चालत नाही हे समजण्यासाठी
  $f(x) \simeq \fn{Un}(x,x)+1$ हे फलन विचारात घ्या. ते आंशिक
  संगणनक्षम आहे का? असल्यास, त्याला निर्देशांक~$e$ आहे, म्हणजे
  $f(x) \simeq \fn{Un}(e,x)$. मग~$f(e)$ विषयी काय म्हणता येईल?
\end{prob}


\end{document}
